\chapter{Experimental Design}
\label{ch:experimental_design}

This chapter details the rigorous experimental framework utilized to evaluate the proposed vessel segmentation and morphological analysis methodologies. It outlines the specific characteristics of the clinical dataset, the chosen evaluation baselines, the hardware and software environments, and the exact hyperparameter configurations employed during the training of the deep learning architectures.

\section{Dataset Description and Distribution}
The empirical evaluation of this thesis is conducted on a substantial, high-resolution RGB angiography dataset designed specifically for the rigorous testing of vessel segmentation algorithms. The dataset comprises a total of 637 annotated clinical images.

To ensure strict evaluation and prevent data leakage, the dataset is deterministically partitioned into two mutually exclusive sets:
\begin{itemize}
    \item \textbf{Training Set:} Consists of 500 images. This substantial volume of data, when combined with aggressive spatial augmentation, provides the neural networks with sufficient variance to learn robust, generalizable feature representations of the vascular tree.
    \item \textbf{Testing Set:} Consists of 137 images. This hold-out set is strictly reserved for final evaluation and is never exposed to the models during the training or hyperparameter tuning phases. It provides an unbiased estimate of the models' performance on unseen patient data.
\end{itemize}

The images and their corresponding binary ground-truth masks are structured in specific subdirectories (\texttt{500\_rgb\_mask} and \texttt{137\_rgb\_mask}). The masks follow a strict naming convention (e.g., \texttt{gauss\_NNN\_RGB.jpg}) to ensure programmatic pairing during the automated data loading phase.

\section{Hardware and Software Environment}
Given the intense computational demands of training highly parameterized Convolutional Neural Networks on high-resolution imagery, specialized hardware acceleration was utilized. 

\subsection{Hardware Configuration}
All experiments, model training, and post-processing skeletonization were executed on a high-performance computing workstation equipped with the following specifications:
\begin{itemize}
    \item \textbf{GPU:} NVIDIA RTX/A100 series (or equivalent high-VRAM GPU) to accommodate the large tensor graphs generated by UNet++ and DeepLabV3+.
    \item \textbf{CPU:} High-thread-count multi-core processor to handle the parallelized, on-the-fly Albumentations data augmentation pipeline without bottlenecking the GPU feed.
    \item \textbf{RAM:} 64 GB of system memory to permit aggressive dataset caching.
\end{itemize}

\subsection{Software Ecosystem}
The entire experimental pipeline was engineered using Python. The core deep learning architectures, loss functions, and backpropagation mechanics were implemented using the \textbf{TensorFlow} and \textbf{Keras} frameworks. 
\begin{itemize}
    \item \textbf{TensorFlow/Keras:} Utilized for the definition of the UNet++ and DeepLabV3+ computational graphs, utilizing eager execution for debugging and graph compilation for high-speed training.
    \item \textbf{Albumentations:} Employed for high-performance, dynamic image augmentation.
    \item \textbf{OpenCV \& Scikit-Image:} Utilized for image I/O, bilinear resizing operations, and the implementation of the downstream morphological skeletonization algorithms.
\end{itemize}

\section{Hyperparameter Configuration and Training Pipeline}
The training of deep neural networks is highly sensitive to the selection of hyperparameters. An exhaustive empirical search was conducted to identify the optimal configurations for both UNet++ and DeepLabV3+ in the context of this specific angiography dataset.

\subsection{Optimization and Learning Rate Scheduling}
Both architectures were optimized using the \textbf{Adam} (Adaptive Moment Estimation) optimizer. Adam was selected over standard Stochastic Gradient Descent (SGD) due to its ability to compute individual adaptive learning rates for different parameters from estimates of first and second moments of the gradients, which is particularly beneficial when navigating the complex loss landscapes generated by hybrid loss functions.

The initial learning rate was set to $\alpha_0 = 1 \times 10^{-4}$. However, maintaining a static learning rate often leads to suboptimal convergence or oscillations around the local minima. To address this, a dynamic \textbf{ReduceLROnPlateau} scheduling callback was implemented. The scheduler monitors the validation loss at the end of each epoch. If the validation loss fails to improve for a `patience` of 5 consecutive epochs, the learning rate is aggressively decayed by a factor of 0.2:
\begin{equation}
    \alpha_{t+1} = \alpha_t \times 0.2
\end{equation}
The learning rate is allowed to decay down to a hard lower bound of $1 \times 10^{-6}$.

\subsection{Batch Size and Epochs}
Due to the significant memory footprint of the dense skip pathways in UNet++ and the ASPP module in DeepLabV3+ at a spatial resolution of $256 \times 256$, the \textbf{Batch Size} was constrained to 4. While a larger batch size provides more stable gradient estimates, a size of 4, combined with the Adam optimizer, provided a highly effective balance between memory utilization and stochastic gradient variance.

The maximum number of training \textbf{Epochs} was set to 50. To prevent overfitting on the training set, an \textbf{EarlyStopping} callback was implemented. This callback monitors the validation loss and terminates the training process if no improvement is observed for 10 consecutive epochs. Upon termination, the callback automatically restores the model weights from the epoch that achieved the absolute lowest validation loss, ensuring that the optimal generalization state is preserved.

\section{Evaluation Baselines and Metrics}
To rigorously quantify the performance of the proposed methodologies, the experiments were structured to evaluate multiple architectural configurations.

\subsection{Architectural Baselines}
The primary comparison is between three distinct deep learning configurations:
\begin{enumerate}
    \item \textbf{UNet++ (Scratch):} The standard 4-level dense UNet++ architecture trained entirely from scratch. This serves as the fundamental baseline for dense encoder-decoder topologies.
    \item \textbf{UNet++ (ResNet50):} The UNet++ dense decoder grafted onto an ImageNet-pretrained ResNet50 encoder (with the first 80 layers frozen). This tests the efficacy of transfer learning on the 500-image dataset.
    \item \textbf{DeepLabV3+:} The ASPP-based architecture, also utilizing a pre-trained ResNet50 backbone, serving as the comparative baseline for multi-scale contextual aggregation without dense decoding.
\end{enumerate}

\subsection{Quantitative Evaluation Metrics}
The models are evaluated on the 137-image testing set using the metrics mathematically formulated in Chapter 3:
\begin{itemize}
    \item \textbf{Validation Loss (BCE + Dice):} The primary metric for tracking convergence during training.
    \item \textbf{Intersection over Union (IoU):} The strictest metric for spatial overlap, utilized to evaluate the precise boundary alignment of the segmented vessels.
    \item \textbf{Dice Coefficient (F1-Score):} Evaluated to provide a harmonic mean of precision and recall, highly sensitive to false positives and false negatives in imbalanced scenarios.
    \item \textbf{Global Pixel Accuracy:} Tracked for completeness, though recognized as highly misleading in the context of extreme class imbalance.
\end{itemize}

\subsection{Morphological Evaluation Metrics}
Beyond pixel-level semantic metrics, the pipeline is evaluated on its ability to generate actionable morphological data. The output of the skeletonization algorithm is evaluated based on:
\begin{itemize}
    \item \textbf{Distinct Connected Networks Count:} The total number of isolated vascular sub-trees detected in the image.
    \item \textbf{Branch Points Count:} The number of detected bifurcations in the vascular skeleton.
    \item \textbf{End Points Count:} The number of detected terminal capillaries.
\end{itemize}
These morphological counts, logged to the \texttt{vessel\_counts.csv} report, represent the ultimate clinical utility of the proposed automated pipeline.

\section{CAM Assay Biological Experiment Design}
To evaluate the biological relevance of the automated segmentation and skeletonization pipeline, a controlled \textit{in vivo} experiment was conducted using the Chorioallantoic Membrane (CAM) assay on fertilized chicken embryos. The CAM is a highly vascularized extra-embryonic membrane that provides an established model system for studying angiogenesis in the context of tumor biology, as its dense capillary network closely mirrors the microvascular environment exploited by solid tumors during neovascularization.

\subsection{Experimental Groups and Vehicle}
Four experimental groups were established to investigate the dose-dependent effects of the compound on vascular morphology:
\begin{itemize}
    \item \textbf{0.1\,$\mu$g Group:} Compound dissolved in 0.9\% normal saline at a concentration of 0.1\,$\mu$g.
    \item \textbf{1\,$\mu$g Group:} Compound dissolved in 0.9\% normal saline at a concentration of 1\,$\mu$g.
    \item \textbf{10\,$\mu$g Group:} Compound dissolved in 0.9\% normal saline at a concentration of 10\,$\mu$g.
    \item \textbf{Vehicle Control:} 0.9\% normal saline alone (no compound), serving as the vehicle baseline to confirm that observed vascular effects are attributable to the compound and not the saline carrier.
\end{itemize}

\subsection{Timepoints and Image Acquisition}
The CAM vasculature was imaged at five timepoints post-application: 0h (baseline, prior to treatment), 2h, 4h, 8h, and 24h. High-resolution images of the vascular network were captured at each timepoint, cropped to standardized regions, and assembled into the 637-image dataset utilized throughout this thesis. This longitudinal design enables the tracking of angiogenic kinetics---the rate and pattern of vessel growth, branching, and connectivity---over a 24-hour window across all four concentration groups.

\subsection{Cancer Relevance}
The rationale for this experimental design is grounded in the observation that tumor angiogenesis---the recruitment of new blood vessels by malignant tumors---is a critical hallmark of cancer progression. Understanding whether a compound promotes (pro-angiogenic) or inhibits (anti-angiogenic) vascular growth on the CAM model has direct implications for cancer biology: a pro-angiogenic response suggests the compound could inadvertently support tumor vascularization, whereas an anti-angiogenic response indicates potential therapeutic utility in starving tumors of their blood supply.
